% ========================================================
% BÀI TẬP TOÁN 5 - KẾT NỐI TRI THỨC VỚI CUỘC SỐNG
% Bài 66: Thực hành và trải nghiệm thu thập, phân tích số liệu
% Định dạng: BTTeX Part III - Trắc nghiệm điền số ngắn
% ========================================================

% Câu 1
\begin{ex}
  Một tổ học sinh khảo sát thời gian tự học ở nhà mỗi ngày của 10 bạn (phút): 60, 90, 45, 60, 120, 90, 60, 60, 90, 45. Thời gian tự học nhiều nhất là bao nhiêu phút?
  \shortans{120}
  \loigiai{
    Giá trị lớn nhất là 120 phút.
  }
\end{ex}

% Câu 2
\begin{ex}
  Ở bảng số liệu trên, thời gian tự học ít nhất là bao nhiêu phút?
  \shortans{45}
  \loigiai{
    Giá trị nhỏ nhất là 45 phút.
  }
\end{ex}

% Câu 3
\begin{ex}
  Ở bảng trên, có bao nhiêu bạn dành đúng 60 phút mỗi ngày để tự học?
  \shortans{4}
  \loigiai{
    Có 4 bạn dành 60 phút.
  }
\end{ex}

% Câu 4
\begin{ex}
  Số lượng cây xanh của 5 lớp khối 5 trồng được là: 5A: 40 cây, 5B: 50 cây, 5C: 45 cây, 5D: 55 cây, 5E: 60 cây. Trung bình mỗi lớp trồng được bao nhiêu cây?
  \shortans{50}
  \loigiai{
    $(40 + 50 + 45 + 55 + 60) : 5 = 250 : 5 = 50$ cây.
  }
\end{ex}

% Câu 5
\begin{ex}
  Ở dự án trên, tổng số cây cả 5 lớp trồng được là bao nhiêu cây?
  \shortans{250}
  \loigiai{
    Tổng là 250 cây.
  }
\end{ex}

% Câu 6
\begin{ex}
  Đo nhịp tim sau vận động của 4 bạn: 110, 120, 115, 135 nhịp/phút. Nhịp tim trung bình của 4 bạn là bao nhiêu nhịp/phút?
  \shortans{120}
  \loigiai{
    $(110 + 120 + 115 + 135) : 4 = 480 : 4 = 120$ nhịp/phút.
  }
\end{ex}

% Câu 7
\begin{ex}
  Thống kê số chai nhựa thu gom trong một tuần của một lớp: Thứ 2: 25 chai, Thứ 3: 30 chai, Thứ 4: 35 chai, Thứ 5: 40 chai, Thứ 6: 70 chai. Cả tuần lớp thu gom được bao nhiêu chai nhựa?
  \shortans{200}
  \loigiai{
    $25 + 30 + 35 + 40 + 70 = 200$ chai.
  }
\end{ex}

% Câu 8
\begin{ex}
  Ở tuần thu gom trên, số chai thu được vào Thứ 6 chiếm bao nhiêu phần trăm cả tuần?
  \shortans{35}
  \loigiai{
    $(70 : 200) \times 100\% = 35\%$.
  }
\end{ex}

% Câu 9
\begin{ex}
  Một hộ gia đình ghi lại số kWh điện dùng trong 3 tháng đầu năm: Tháng 1: 180 kWh, Tháng 2: 170 kWh, Tháng 3: 250 kWh. Trung bình mỗi tháng gia đình dùng bao nhiêu kWh?
  \shortans{200}
  \loigiai{
    $(180 + 170 + 250) : 3 = 600 : 3 = 200\text{ kWh}$.
  }
\end{ex}

% Câu 10
\begin{ex}
  Một nhóm học sinh đo nhiệt độ trong phòng lúc 12 giờ trưa trong 5 ngày liên tiếp: $28^\circ\text{C}, 30^\circ\text{C}, 31^\circ\text{C}, 29^\circ\text{C}, 32^\circ\text{C}$. Nhiệt độ trung bình là bao nhiêu độ C?
  \shortans{30}
  \loigiai{
    $(28 + 30 + 31 + 29 + 32) : 5 = 150 : 5 = 30^\circ\text{C}$.
  }
\end{ex}

% Câu 11
\begin{ex}
  Số lượng học sinh cận thị của khối 5 là 40 bạn trên tổng số 160 bạn. Tỉ số phần trăm học sinh cận thị là bao nhiêu phần trăm?
  \shortans{25}
  \loigiai{
    $(40 : 160) \times 100\% = 25\%$.
  }
\end{ex}

% Câu 12
\begin{ex}
  Ở khối 5 trên, có bao nhiêu bạn không bị cận thị?
  \shortans{120}
  \loigiai{
    $160 - 40 = 120$ bạn.
  }
\end{ex}

% Câu 13
\begin{ex}
  Thực nghiệm gieo xúc xắc 100 lần, ghi nhận được kết quả số chấm chẵn là 54 lần. Số lần xuất hiện số chấm lẻ là bao nhiêu lần?
  \shortans{46}
  \loigiai{
    $100 - 54 = 46$ lần.
  }
\end{ex}

% Câu 14
\begin{ex}
  Một khảo sát về loại quả yêu thích của 50 bạn: Cam 20 bạn, Táo 15 bạn, Xoài 10 bạn, Chuối 5 bạn. Tỉ số phần trăm số bạn thích cam là bao nhiêu phần trăm?
  \shortans{40}
  \loigiai{
    $(20 : 50) \times 100\% = 40\%$.
  }
\end{ex}

% Câu 15
\begin{ex}
  Ở khảo sát trên, tỉ số phần trăm số bạn thích chuối là bao nhiêu phần trăm?
  \shortans{10}
  \loigiai{
    $(5 : 50) \times 100\% = 10\%$.
  }
\end{ex}

% Câu 16
\begin{ex}
  Đo chiều dài bước chân của 4 học sinh (cm): 50, 52, 48, 50. Chiều dài bước chân trung bình là bao nhiêu cm?
  \shortans{50}
  \loigiai{
    $(50 + 52 + 48 + 50) : 4 = 200 : 4 = 50\text{ cm}$.
  }
\end{ex}

% Câu 17
\begin{ex}
  Bảng theo dõi số trang sách Nam đọc mỗi ngày trong một tuần: 10, 15, 20, 15, 20, 25, 35. Cả tuần Nam đọc được bao nhiêu trang sách?
  \shortans{140}
  \loigiai{
    $10 + 15 + 20 + 15 + 20 + 25 + 35 = 140$ trang.
  }
\end{ex}

% Câu 18
\begin{ex}
  Ở bảng trên, trung bình mỗi ngày Nam đọc được bao nhiêu trang sách?
  \shortans{20}
  \loigiai{
    $140 : 7 = 20$ trang.
  }
\end{ex}

% Câu 19
\begin{ex}
  Một nông dân ghi chép số trứng gà thu hoạch trong 4 ngày: 120, 115, 130, 135 quả. Tổng số trứng thu hoạch được là bao nhiêu quả?
  \shortans{500}
  \loigiai{
    $120 + 115 + 130 + 135 = 500$ quả.
  }
\end{ex}

% Câu 20
\begin{ex}
  Ở trang trại trên, trung bình mỗi ngày thu được bao nhiêu quả trứng?
  \shortans{125}
  \loigiai{
    $500 : 4 = 125$ quả.
  }
\end{ex}

% Câu 21
\begin{ex}
  Khảo sát thời gian xem ti-vi của 20 bạn: 5 bạn xem dưới 30 phút, 10 bạn xem từ 30 đến 60 phút, 5 bạn xem trên 60 phút. Tỉ số phần trăm bạn xem trên 60 phút là bao nhiêu phần trăm?
  \shortans{25}
  \loigiai{
    $(5 : 20) \times 100\% = 25\%$.
  }
\end{ex}

% Câu 22
\begin{ex}
  Thực hành đo thể tích nước tiêu thụ của gia đình trong 4 tuần: $4\text{ m}^3, 5\text{ m}^3, 6\text{ m}^3, 5\text{ m}^3$. Tổng lượng nước tiêu thụ là bao nhiêu mét khối?
  \shortans{20}
  \loigiai{
    $4 + 5 + 6 + 5 = 20\text{ m}^3$.
  }
\end{ex}

% Câu 23
\begin{ex}
  Nếu giá mỗi mét khối nước ở câu trên là $10\,000$ đồng, tổng tiền nước trong 4 tuần là bao nhiêu nghìn đồng?
  \shortans{200}
  \loigiai{
    $20 \times 10 = 200$ nghìn đồng.
  }
\end{ex}

% Câu 24
\begin{ex}
  Thống kê số điểm 10 của lớp trong 4 tháng: Tháng 9: 45 điểm 10, Tháng 10: 60, Tháng 11: 55, Tháng 12: 80. Tổng số điểm 10 cả học kì I là bao nhiêu?
  \shortans{240}
  \loigiai{
    $45 + 60 + 55 + 80 = 240$.
  }
\end{ex}

% Câu 25
\begin{ex}
  Ở lớp học trên, trung bình mỗi tháng có bao nhiêu điểm 10?
  \shortans{60}
  \loigiai{
    $240 : 4 = 60$ điểm 10.
  }
\end{ex}
